\documentclass[10pt,a4paper]{amsart}
\usepackage[margin=19mm]{geometry}
\usepackage{amsmath,amssymb,mathtools}
\usepackage[scr=rsfs,cal=euler]{mathalfa}
\usepackage{xfrac}
\usepackage[unicode,hidelinks,bookmarks=false]{hyperref}
\newcommand{\ie}{i.e.\ }
\newcommand{\cf}{cf.\ }
\newcommand{\resp}{respectively\ }
\newcommand{\Subsection}{\subsection}
\providecommand{\nobreakdash}{\nobreak\hbox{-}}
\setlength{\emergencystretch}{2em}
\multlinegap0pt
\usepackage{bm}
\usepackage{bbm}
\usepackage{dsfont}
\usepackage{xspace}
\usepackage{cancel}
\usepackage{mathtools}
\usepackage{amsthm}
\makeatletter
\@ifundefined{thm}{\newtheorem{thm}{Theorem}}{}
\@ifundefined{lem}{\newtheorem{lem}[thm]{Lemma}}{}
\@ifundefined{prop}{\newtheorem{prop}[thm]{Proposition}}{}
\@ifundefined{rem}{\newtheorem{rem}[thm]{Remark}}{}
\makeatother
\makeatletter
\newcommand{\Ncal}{\mathcal{N}}
\def\dfrac{\displaystyle\frac}
\def\diam{{\mathop\mathrm{\,diam\,}}}
\def\dist{{\mathop\mathrm{dist}}}
\def\ez{\epsilon}
\@ifundefined{note}{\newenvironment{note}{\begin{quote}\sf}{\end{quote}}}{}
\makeatother
\title[An anisotropic Serrin's problem in general domains]{An anisotropic Serrin's problem in general domains}
\author{Alessio Figalli, Yi Ru-Ya Zhang}
\date{Source: arXiv:2603.06119v2 (CC BY 4.0)}
\begin{document}
\maketitle

\noindent\emph{Source and licence.} An anisotropic Serrin's problem in general domains. Version arXiv:2603.06119v2 (CC BY 4.0), distributed under the Creative Commons Attribution 4.0 International licence, as recorded in the arXiv licence metadata for this exact version and shown on the abstract page https://arxiv.org/abs/2603.06119v2. Full rights notice: licenses/arxiv-2603.06119-CC-BY-4.0.txt.\par\smallskip

\noindent\textit{Editorial scope.} Counted unit: the Lem whose source label is \texttt{second derivative} in the article (source0.tex lines 708--719, source label \texttt{second derivative}), delivered together with the article's own complete original proof of it (source0.tex lines 721--1078, one \texttt{proof} environment of 358 lines). No mathematics is written, repaired or re-derived: statement and proof are the article's own text line for line. Required context retained uncounted: main thm (lines 369--392); beta number (lines 371--373); domain boundary1 (lines 375--377); eq:weak set finite per (lines 381--384); eq:minkowski-content (lines 402--413); basic property u (lines 460--501); prop:ADR (lines 1738--1777); item:badcount (lines 1759--1776); lem:flat-cover (lines 1888--1924). Every definition, display and label of those lines is retained unchanged. Omitted from this package and not redistributed: 1--368, 374--374, 378--380, 385--401, 414--459, 502--707, 720--720, 1079--1737, 1777--1887, 1925--2121. Nothing inside a retained excerpt is omitted or altered: every retained body is delivered exactly as the article's lines are written, so no omission receipt of any kind accompanies this package. Citations: the retained excerpts carry no citation command at all, so no bibliography entry of the article is transcribed and no reference is redistributed. Reference adaptation: none was needed and none was made; every reference inside the delivered unit resolves inside it, so no unresolved reference or citation is produced. Printed slips kept literal: no printed word, symbol, hypothesis or number of the counted statement or of its proof was altered. Layout and class: the article's own class is \texttt{amsart} with packages including a4wide, color, eucal, enumerate, mathrsfs, ulem, amsmath, amssymb, epsfig, amsthm, inputenc, psfrag, hyperref, cleveref; the wrapper is the lane's pinned \texttt{amsart} editorial wrapper at 10pt on A4, which restates the theorem-like environments the retained text needs and adds the article's own macro definitions that the retained text needs (\texttt{\textbackslash Ncal}, \texttt{\textbackslash dfrac}, \texttt{\textbackslash diam}, \texttt{\textbackslash dist}, \texttt{\textbackslash ez}), with extra wrapper packages (bm, bbm, dsfont, xspace, cancel, mathtools). The delivered numbering is the unit's own. Assets: no retained span contains a figure environment, a graphics-inclusion command, a PDF-page inclusion, a table or a TikZ picture; no image file is copied into this package, the package declares no asset dependency and no image file is redistributed with it. Licence: version arXiv:2603.06119v2 of this article is distributed under the Creative Commons Attribution 4.0 International licence, as recorded in the arXiv licence metadata for this exact version and shown on the abstract page https://arxiv.org/abs/2603.06119v2. A full rights notice is delivered at licenses/arxiv-2603.06119-CC-BY-4.0.txt. Characters: every retained excerpt is pure ASCII, so the delivered engine, XeLaTeX with its default Latin Modern text font, reports no missing character.
\begin{thm}\label{main thm}
Let $\Omega\subset \mathbb R^n$ be a bounded indecomposable set of finite perimeter,  and assume there exist constant $A_1,A_2\geq 1$ such that 
\begin{equation}\label{beta number}
\int_{\partial^* \Omega}\int_{0}^{1} \beta(x,s)^2 \, \frac{ds} {s} \, d\mathscr H^{n-1}(x)\le A_1
\end{equation}
and that, for every $x\in \partial^* \Omega$ and every $r \in (0,1)$,
\begin{equation}\label{domain boundary1}
A_2^{-1}r^{n-1}\le \mathscr H^{n-1}(B_r(x)\cap \partial^* \Omega)\le A_2 r^{n-1}.
\end{equation}
Let $K$ be a bounded uniformly convex body whose boundary is of class
$C^{2,\gamma}$ for some $\gamma>0$ small, and let $H$ be the corresponding Wulff potential.
Then $\Omega$ admits a solution $u\in W^{1,2}(\mathbb R^n)$ to
\begin{equation}\label{eq:weak set finite per}
u=0\quad \text{a.e.\ in }\mathbb R^n\setminus \Omega,\qquad
\Delta_H u=\mathbf{c}H(-\nu)\,\mathscr{H}^{n-1}\llcorner \partial^*\Omega - \mathbf{1}_{\Omega}\,dx,
\end{equation}
in the sense of distributions, if and only if, {up to modifying $\Omega$ on a null set and up to a translation,} 
\[
{
\text{$\Omega$ is homothetic to $-K$}\qquad \text{and}\qquad
u(x)=\dfrac{(r^2-H_*^2(-x))_+}{2n}\quad \text{for some }r>0.
}
\]
\end{thm}

\begin{rem}
We say that a set $\Omega\subset\mathbb R^n$ has finite $(n-1)$-dimensional upper Minkowski content if there exists $M>0$ such that
\begin{equation}\label{eq:minkowski-content}
    {\limsup_{r\to 0^+} \frac{|\mathcal N_{r}(\partial\Omega)|}{r}\le M<\infty,}
\end{equation}
where, for $r>0$, the $r$-neighborhood of $\partial\Omega$ is
\[
\mathcal N_r(\partial\Omega):=\left\{ x\in \mathbb R^n:\ \dist(x,\partial \Omega)\le r\,\diam(\Omega) \right\}.
\]
Note that, via a Besicovitch covering, one can directly conclude from \eqref{domain boundary1} that $\partial \Omega$ has a finite
Minkowski content with $M=M(n,A_2)$.
\end{rem}

\begin{lem}\label{basic property u}
Let $\Omega\subset \mathbb R^n$ be a bounded set of finite perimeter satisfying \eqref{domain boundary1},
and let $u \in W^{1,2}(\mathbb R^n)$ satisfy \eqref{eq:weak set finite per}, where $H$ satisfies the assumptions in Theorem~\ref{main thm}.
Then:
\begin{enumerate}
\item[(1)] $u$ is nonnegative and $\mathring\Omega=\{u>0\}$.

\item[(2)] $u$ is {Lipschitz continuous} (globally in $\mathbb R^n$).

\item[(3)] $\Omega=\mathring\Omega$ up to a set of measure zero. In particular, without loss of generality, we may assume
$\Omega$ to be open.

\item[(4)] For every $x\in \partial^*\Omega$ it holds
\[
\frac{u(x+rz)}{r} \to  \mathbf{a}(x)\bigl(-\nu_x\cdot z \bigr)_+ \quad \text{ as } \ r\to 0,
\]
where $\nu_x$ denotes the measure-theoretic outer unit normal at $x$ and $\mathbf {a}(x):=\frac{\mathbf c}{H(-\nu_x)}$.
{Moreover, the interior gradients have the approximate trace
\[
\nabla u(y)\to -\mathbf a(x)\nu_x
\quad\text{as } y\to x \text{ from inside }\Omega
\]
in the measure-theoretic sense, i.e. after rescaling at $x$ the gradients converge in measure on interior half-balls.}

\item[(5)] For any $0<\kappa<1$, any $\delta>0$, any $\eta>0$, and every $y\in \partial^*\Omega$, there exists
$r_0=r_0(y,n,\kappa,\delta,\eta)>0$ such that for every $0<r<r_0$ one has
\[
\beta(y,r)\le \eta,
\qquad
\frac{\mathscr {H}^{n-1}(\partial^*\Omega\cap B(y,r))}{\omega_{n-1}r^{n-1}}\in (1-\eta,\,1+\eta),
\]
and
\[
\int_{\hat B_{y,\,\kappa,\, r}}|D^2u|\, dx\le \delta \,(\kappa^{-1} r)^{n-1},
\]
where
\[
\hat B_{y,\,\kappa,\, r}:=
B_{\frac{\kappa^{-1}-\kappa}{2}r}\!\left(y-\frac{\kappa^{-1}+\kappa}{2}r\,\nu_y\right){\subset\subset \Omega}.
\]
\end{enumerate}
\end{lem}

\begin{prop}[ADR covering with density control and $1/|\log r|$ gain on $\beta$-bad]\label{prop:ADR}
Fix $\eta\in(0,1)$ and let $\theta:(0,1)\to(0,1)$ be any function with $\theta(\eta)\downarrow 0$ as $\eta\downarrow 0$.
Then there exist $r_0=r_0(\Omega,\eta,{\rm ADR})\in(0,r_{\rm ADR})$ and a Borel set $F_\eta\subset\partial^*\Omega$ such that, for all $x \in F_\eta$ and $0<s\leq r_0$,
\begin{equation}\label{eq:Egorov}
  \mathscr H^{n-1}(\partial^*\Omega\setminus F_\eta)\ \le\ \theta(\eta)\,\eta^{\,n-1},\qquad
  \frac{\mathscr H^{n-1}(\partial^*\Omega\cap B(x,s))}{\omega_{n-1}s^{n-1}}\in\Big(1-\tfrac{\eta}{10},\,1+\tfrac{\eta}{10}\Big).
\end{equation}
{Moreover, let $G_\eta\subset\partial^*\Omega$ be any Borel set satisfying}
\[
{
\mathscr H^{n-1}(\partial^*\Omega\setminus G_\eta)\le \theta(\eta)\,\eta^{n-1}.
}
\]
{Then there exists a sequence $r_j\downarrow0$ such that, for each $j$, there is a Besicovitch covering}
\[
{
\mathcal B_{r_j}=\mathcal G_{r_j}\cup\mathcal R_{r_j}
}
\]
{of $\Ncal_{r_j/2}(\partial\Omega)$ by balls $B(x_h,r_j)$ with centers in $\partial^*\Omega$ and bounded overlap, with the following properties:}

\begin{enumerate}
\item\label{item:gooddensity}
{If $B(x_h,r_j)\in\mathcal G_{r_j}$, then $x_h\in F_\eta\cap G_\eta$, $\beta(x_h,r_j)\le\eta$, and}
\[
{
  \frac{\mathscr H^{n-1}(\partial^*\Omega\cap B(x_h,r_j))}{\omega_{n-1}r_j^{n-1}}\in(1-\eta,\,1+\eta).
}
\]

\item\label{item:badcount}
{The bad subfamily satisfies}
\[
{
  \#\mathcal R_{r_j}\ \le\ C(n,{\rm ADR})\,r_j^{-(n-1)}
  \left(\frac{\eta^{-2}A_1}{|\log r_j|}+\theta(\eta)\right).
}
\]
\end{enumerate}
\end{prop}

\begin{lem}[Geometric covering of the intermediate layer]\label{lem:flat-cover}
Fix $\kappa\in(0,1)$. There exist $\bar\kappa=\bar\kappa(\kappa)\in(0,\kappa)$, $\eta_0=\eta_0(n,\kappa)>0$, and
$C_\kappa=C(n,\kappa)$ with the following property. Set
\[
c_\kappa:=(4+\kappa^{-1}+\bar\kappa^{-1})^{-1}.
\]
Let $0<\eta\le\eta_0$, let $\theta$ be the function used in Proposition~\ref{prop:ADR}, and let
$G_{\kappa,\eta}\subset\partial^*\Omega$ satisfy
\[
\mathscr H^{n-1}(\partial^*\Omega\setminus G_{\kappa,\eta})\le \theta(\eta)\,\eta^{n-1}.
\]
Let $r_{\kappa,\eta}>0$ be such that, for every
$y\in G_{\kappa,\eta}$ and every $0<s\le r_{\kappa,\eta}$,
\[
\beta(y,s)\le\eta,\qquad 
\frac{\mathscr H^{n-1}(\partial^*\Omega\cap B(y,s))}{\omega_{n-1}s^{n-1}}\in(1-\eta,1+\eta).
\]
Let $0<r<r_{\rm ADR}$ with $r\le c_\kappa r_{\kappa,\eta}$, and let
$\{B(y_h,r)\}_h=\mathcal G_r\cup\mathcal R_r$ be a covering given by Proposition~\ref{prop:ADR} with 
$G_\eta=G_{\kappa,\eta}$.
After the normalization $0\in\Omega$ and $\diam(\Omega)=1$, set
\[
E_{\kappa,r}:=\{x\in\Omega:\ \kappa r<\dist(x,\partial\Omega)\le \kappa^{-1}r,\quad
(1+r)x\in\Omega\setminus\mathcal N_{\kappa r}(\partial\Omega)\}.
\]
For $B(y_h,r)\in\mathcal G_r$ set
\[
\widehat B_h:=\widehat B_{y_h,\bar\kappa,r},\qquad
E_h:=\{x\in E_{\kappa,r}:\ x\in \widehat B_h,\ (1+r)x\in \widehat B_h\}.
\]
Then
\[
E_{\kappa,r}\setminus\bigcup_{B(y_h,r)\in\mathcal G_r}E_h
\subset
\bigcup_{B(y_h,r)\in\mathcal R_r} B(y_h,C_\kappa r).
\]
\end{lem}

\begin{lem}\label{second derivative}
Let $\Omega$ be a bounded set of finite perimeter satisfying \eqref{beta number} and \eqref{domain boundary1}, and let
$u$ be a weak solution of \eqref{eq:weak set finite per}. Then
\begin{multline}\label{change derivative}
\liminf_{\varepsilon\to 0}\bigg|
\int_{\Omega} DV(\nabla u(x))\cdot \frac{\alpha_{\varepsilon}(x)-\alpha_{-\varepsilon}(x)}{2\varepsilon}\,dx \\
-
\int_{\Omega} \nabla u(x)\cdot \frac{DV(\alpha_{\varepsilon}(x))-DV(\alpha_{-\varepsilon}(x))}{2\varepsilon}\,dx
\bigg|=0,
\end{multline}
where $\alpha_{\pm\varepsilon}(x):=\nabla u((1\pm\varepsilon)x)$.
\end{lem}

\begin{proof}
Recall that $u\in W^{2,2}_{\rm loc}(\Omega)$ and that $V(\xi)=\frac12 H(\xi)^2$ is $2$-homogeneous. Hence, by Euler's theorem,
\[
DV(\xi)=D^2V(\xi)\,\xi \qquad \text{for }\xi\neq 0.
\]
Up to a translation (fixed once and for all), we may assume that $0\in\Omega$. In particular,
\begin{equation}\label{eq:xbound-diam}
|x|\le \diam(\Omega)\qquad \forall\,x\in\Omega.
\end{equation}
{Since the assertion is invariant under dilations of the independent variable, we also normalize
$\diam(\Omega)=1$ throughout the proof.}
We also recall that $u=0$ a.e.\ on $\mathbb R^n\setminus\Omega$, hence $\nabla u=0$ a.e.\ on $\mathbb R^n\setminus\Omega$.

\medskip
\noindent\textit{Step 1: Algebraic reduction to two symmetric terms.}
Set
\[
A_\varepsilon:=
\int_{\Omega} DV(\nabla u)\cdot \frac{\alpha_{\varepsilon}-\alpha_{-\varepsilon}}{2\varepsilon}\,dx
-
\int_{\Omega} \nabla u\cdot \frac{DV(\alpha_{\varepsilon})-DV(\alpha_{-\varepsilon})}{2\varepsilon}\,dx.
\]
Then
\begin{align*}
A_\varepsilon
&=\int_{\Omega} \Bigg[
DV(\nabla u)\cdot \frac{\alpha_{\varepsilon}-\nabla u}{2\varepsilon}
-\nabla u\cdot \frac{DV(\alpha_{\varepsilon})-DV(\nabla u)}{2\varepsilon}
\Bigg]dx\\
&\quad +\int_{\Omega} \Bigg[
DV(\nabla u)\cdot \frac{\nabla u-\alpha_{-\varepsilon}}{2\varepsilon}
-\nabla u\cdot \frac{DV(\nabla u)-DV(\alpha_{-\varepsilon})}{2\varepsilon}
\Bigg]dx.
\end{align*}
Define
\begin{align*}
I_\varepsilon^\pm
:=\Bigg|\int_{\Omega} \Bigg[
DV(\nabla u(x))\cdot \frac{\alpha_{\pm\varepsilon}(x)-\nabla u(x)}{2\varepsilon}
-\nabla u(x)\cdot \frac{DV(\alpha_{\pm\varepsilon}(x))-DV(\nabla u(x))}{2\varepsilon}
\Bigg]dx\Bigg|.
\end{align*}
Since the second bracket above is the negative of the integrand defining $I_\varepsilon^-$, we have
\[
|A_\varepsilon|\le I_\varepsilon^+ + I_\varepsilon^-.
\]
Therefore, to prove \eqref{change derivative} it suffices to show
\[
\liminf_{\varepsilon\to 0}\big(I_\varepsilon^+ + I_\varepsilon^-\big)=0.
\]
We begin by analyzing $I_\varepsilon^+$.

\medskip
\noindent\textit{Step 2: Restriction to the set $U$ and Taylor expansion.}
Note that
$$\alpha_\varepsilon(x)=\nabla u((1+\varepsilon)x)=0  \quad \text{ for a.e.  } \ x\in \Omega \ \text{ while } \ (1+\ez)x\notin \Omega,$$ 
and one checks directly that
the integrand in $I_\varepsilon^+$ vanishes (the two terms cancel, using $DV(0)=0$).
Thus we may restrict the integration to
\[
U:=\{x\in\Omega:\ (1+\varepsilon)x\in\Omega\}.
\]
On $U$, by the fundamental theorem of calculus,
\[
DV(\alpha_\varepsilon)-DV(\nabla u)
=\int_0^1 D^2V\bigl(\nabla u+t(\alpha_\varepsilon-\nabla u)\bigr)\,(\alpha_\varepsilon-\nabla u)\,dt,
\]
and using $DV(\xi)=D^2V(\xi)\xi$ we obtain
\begin{align*}
I_\varepsilon^+
&=\Bigg|\int_{U}\nabla u(x)\cdot \Bigg[
D^2V(\nabla u(x))
-\int_0^1 D^2V\bigl(\nabla u(x)+t(\alpha_\varepsilon(x)-\nabla u(x))\bigr)\,dt
\Bigg]\frac{\alpha_\varepsilon(x)-\nabla u(x)}{2\varepsilon}\,dx\Bigg|.
\end{align*}

\medskip
\noindent\textit{Step 3: Using the $C^\gamma$ regularity of $D^2V$ on the sphere.}
Let $\gamma_0 \in (0,1]$ be such that $D^2V\in C^{\gamma_0}(\mathbb S^{n-1})$, and fix $\gamma \in (0,\gamma_0)$.
Since $D^2V$ is $0$-homogeneous and $C^\gamma$ on $\mathbb S^{n-1}$, for every $\xi\neq 0$ and every $\zeta\in\mathbb R^n$ we have
\[
|D^2V(\xi)-D^2V(\xi+\zeta)|
\le C\,\bigg(\frac{|\zeta|}{|\xi|}\bigg)^\gamma.
\]
Applying this with $\xi=\nabla u(x)$ and $\zeta=t(\alpha_\varepsilon(x)-\nabla u(x))$, and using $|\nabla u|\le L$ a.e.
(Lemma~\ref{basic property u}(2)), we obtain
\begin{align*}
I_\varepsilon^+
&\le C\int_U |\nabla u(x)|\,\frac{|\alpha_\varepsilon(x)-\nabla u(x)|}{2\varepsilon}\,
\int_0^1 \bigg(\frac{t|\alpha_\varepsilon(x)-\nabla u(x)|}{|\nabla u(x)|}\bigg)^\gamma dt\,dx\\
&\le C(L,\gamma)\int_U \frac{|\alpha_\varepsilon(x)-\nabla u(x)|^{1+\gamma}}{2\varepsilon}\,dx.
\end{align*}
Thus it remains to show that
\begin{equation}\label{eq:goal-alpha}
\int_U \frac{|\nabla u((1+\varepsilon)x)-\nabla u(x)|^{1+\gamma}}{\varepsilon}\,dx\to 0
\qquad\text{along a sequence }\varepsilon\downarrow 0.
\end{equation}

\medskip
\noindent\textit{Step 4: Decomposition into $I_{1,\varepsilon}^++I_{2,\varepsilon}^++I_{3,\varepsilon}^+$.}
Recall that
\[
\mathcal N_{\delta}(\partial \Omega):=\left\{ x\in \mathbb R^n:\ \dist(x,\partial \Omega)\le \delta\,\diam(\Omega) \right\}.
\]
For $t>0$ define
\[
U_t:=\Big\{x\in \Omega\setminus \mathcal N_{t}(\partial\Omega):\ (1+\varepsilon)x\in\Omega\Big\}.
\]
Fix $\kappa\in(0,1)$ (to be sent to $0$ later). We decompose
\begin{align*}
\int_U \frac{|\nabla u((1+\varepsilon)x)-\nabla u(x)|^{1+\gamma}}{2\varepsilon}\,dx
&=\int_{U_{\kappa^{-1}\varepsilon}} \frac{|\nabla u((1+\varepsilon)x)-\nabla u(x)|^{1+\gamma}}{2\varepsilon}\,dx\\
&\quad +\int_{U_{\kappa\varepsilon}\setminus U_{\kappa^{-1}\varepsilon}} \frac{|\nabla u((1+\varepsilon)x)-\nabla u(x)|^{1+\gamma}}{2\varepsilon}\,dx\\
&\quad +\int_{U\setminus U_{\kappa\varepsilon}} \frac{|\nabla u((1+\varepsilon)x)-\nabla u(x)|^{1+\gamma}}{2\varepsilon}\,dx\\
&=: I_{1,\varepsilon}^+ + I_{2,\varepsilon}^+ + I_{3,\varepsilon}^+.
\end{align*}

\medskip
\noindent\textit{Step 5: Estimate of $I_{1,\varepsilon}^+$ (interior region).}
On $U_{\kappa^{-1}\varepsilon}$ we have $\dist(x,\partial\Omega)\ge \kappa^{-1}\varepsilon\,\diam(\Omega)$, hence both $x$ and $(1+\varepsilon)x$ {remain at a strictly positive distance from the boundary}.
Since $u\in W^{2,2}_{\rm loc}(\Omega)$, interior $W^{2,2}$ estimates for the uniformly elliptic equation
${\rm div}(DV(\nabla u))=-1$ imply that for each ball $B_{2r}(z)\subset\Omega$,
\[
\|D^2u\|_{L^2(B_r(z))}\le C\,r^{\frac n2-1}.
\]
By H\"older's inequality (since $1+\gamma<2$), this yields
\begin{equation}\label{eq:L1+g-Hess}
\int_{B_r(z)} |D^2u|^{1+\gamma}\,dx \le C\,r^{n-1-\gamma}\qquad\text{whenever }B_{2r}(z)\subset\Omega.
\end{equation}
Moreover, by the fundamental theorem of calculus along $t\mapsto (1+t)x$ and \eqref{eq:xbound-diam},
\[
\nabla u((1+\varepsilon)x)-\nabla u(x)=\int_0^\varepsilon D^2u((1+t)x)\,x\,dt,
\]
hence
\[
\frac{|\nabla u((1+\varepsilon)x)-\nabla u(x)|^{1+\gamma}}{\varepsilon}
\le C\,\varepsilon^\gamma\int_0^1 |D^2u((1+t\varepsilon)x)|^{1+\gamma}\,dt.
\]
Integrating over $U_{\kappa^{-1}\varepsilon}$ and using a Besicovitch covering argument with balls
$B_{r(x)}(x)$ where $r(x)\sim \dist(x,\partial\Omega)$, together with \eqref{eq:L1+g-Hess}, yields\footnote{Here we note that $\partial\Omega$ has finite $(n{-}1)$-dimensional upper Minkowski content (a consequence of
	\eqref{domain boundary1}), {which together with the bounded overlap in the Besicovitch covering gives the counting estimate}
	\[
	{\#\{B_{r_k}(x_k):\,2^{-m}\le r_k\le 2^{-m+1}\}\le C(n,M)\,2^{m(n-1)},}
	\] 
	{therefore}
	$$
	{\sum_k r_k^{n-1-\gamma} \leq C (\kappa^{-1}\varepsilon)^{-\gamma}.}
	$$
} 
\[
I_{1,\ez}^+ \le C(L,n,\diam(\Omega))\,\varepsilon^\gamma \sum_k r_k^{n-1-\gamma}
\le C(L,M,n,\diam(\Omega))\,\kappa^\gamma{,}
\]
{where $M$ is the constant corresponding to the Minkowski content \eqref{eq:minkowski-content}.} 

\medskip
\noindent\textit{Step 6: Estimate of $I_{3,\varepsilon}^+$ (very near the boundary).}
Since $|\nabla u|\le L$ a.e., we have
\[
I_{3,\varepsilon}^+\le C(L)\,\varepsilon^{-1}\,|\mathcal N_{\kappa\varepsilon}(\partial\Omega)|
=O(\kappa),
\]
using again the finite upper Minkowski content: $|\mathcal N_{\rho}(\partial\Omega)|=O(\rho)$ as $\rho\to 0$.


\medskip
\noindent\textit{Step 7: Estimate of $I^+_{2,\varepsilon}$ (intermediate region).}
We first split according to whether $(1+\varepsilon)x$ falls within the thin boundary layer.
By the finite upper Minkowski content, we have
\[
\bigl|\{x\in U_{\kappa\varepsilon}\setminus U_{\kappa^{-1}\varepsilon}:\ (1+\varepsilon)x\in \mathcal N_{\kappa\varepsilon}(\partial\Omega)\}\bigr|
=O(|\mathcal N_{\kappa\varepsilon}(\partial\Omega)|)=O(\kappa\varepsilon),
\]
hence, using $|\nabla u|\le L$,
\begin{equation}\label{2.7}
\int_{\{x\in U_{\kappa\varepsilon}\setminus U_{\kappa^{-1}\varepsilon}:\ (1+\varepsilon)x\in \mathcal N_{\kappa\varepsilon}(\partial\Omega)\}}
\frac{|\nabla u((1+\varepsilon)x)-\nabla u(x)|^{1+\gamma}}{2\varepsilon}\,dx
\le O(\kappa).
\end{equation}
We now consider the remaining points
\[
E:=\Bigl(U_{\kappa\varepsilon}\setminus U_{\kappa^{-1}\varepsilon}\Bigr)\cap\Bigl\{x:\ (1+\varepsilon)x\notin \mathcal N_{\kappa\varepsilon}(\partial\Omega)\Bigr\}.
\]

\smallskip
\noindent\emph{Uniform Hessian smallness on a large subset of $\partial^*\Omega$.}
{Fix $\kappa\in(0,1)$. Let $\bar\kappa=\bar\kappa(\kappa)\in(0,\kappa)$ be the geometric constant provided by
Lemma~\ref{lem:flat-cover}, and choose $\delta_\kappa>0$ so small that}
\begin{equation}\label{eq:deltakappa}
{
C(n,\gamma,A_2)\,\delta_\kappa^{1+\gamma}\,\bar\kappa^{-n}\le \kappa^\gamma .
}
\end{equation}
{Fix $\eta\in(0,1)$ and apply Lemma~\ref{basic property u}(5) with the parameters $\bar\kappa$ and
$\delta_\kappa$. For $\mathscr H^{n-1}$-a.e.\ $y\in\partial^*\Omega$ there exists $r_0(y)>0$ such that for all
$0<r<r_0(y)$,}
\begin{equation}\label{three inequ}
{
\beta(y,r)\le\eta,\qquad 
\frac{\mathscr H^{n-1}(\partial^*\Omega\cap B(y,r))}{\omega_{n-1}r^{n-1}}\in(1-\eta,1+\eta),
\qquad
\int_{\widehat B_{y,\bar\kappa,r}}|D^2u|\,dx \le \delta_\kappa\,(\bar\kappa^{-1}r)^{n-1},
}
\end{equation}
{where $\widehat B_{y,\bar\kappa,r}$ denotes the ball $\hat B_{y,\bar\kappa,r}$ of Lemma~\ref{basic property u}(5).}

{Fix once and for all a function $\theta$ as in Proposition~\ref{prop:ADR}, for instance $\theta(\eta)=\eta$.}
For $m\in\mathbb N$ define
\[
G_m:=\Bigl\{y\in\partial^*\Omega:\ \text{the three inequalities in \eqref{three inequ} hold for all }0<r\le 2^{-m}\Bigr\}.
\]
Since $G_m\uparrow \partial^*\Omega$ up to a null set, we may choose $m=m(\kappa,\eta)$ so large that
\[
\mathscr H^{n-1}(\partial^*\Omega\setminus G_m)\le \theta(\eta)\,\eta^{n-1}.
\]
Set
\[
G_{\kappa,\eta}:=G_m,\qquad r_{\kappa,\eta}:=2^{-m}.
\]
In particular, for every $y\in G_{\kappa,\eta}$ and every $0<r\le r_{\kappa,\eta}$ we have
\begin{equation}\label{7.1}
{
\int_{\widehat B_{y,\bar\kappa,r}}|D^2u|\,dx \le \delta_\kappa\,(\bar\kappa^{-1}r)^{n-1}.
}
\end{equation}

\smallskip
\noindent\emph{Good balls and the estimate on the good part.}
{We choose $\eta\le \eta_0(n,\kappa)$, where $\eta_0$ is given by Lemma~\ref{lem:flat-cover}, and apply
Proposition~\ref{prop:ADR} with the auxiliary set $G_\eta=G_{\kappa,\eta}$. Let $\{r_j\}$ and the coverings
$\mathcal B_{r_j}=\mathcal G_{r_j}\cup\mathcal R_{r_j}$ be the resulting sequence and good/bad decomposition.
Let $c_\kappa$ be the constant from Lemma~\ref{lem:flat-cover}. In the rest of the proof we restrict to
$\varepsilon=r_j\le c_\kappa r_{\kappa,\eta}$; this only discards finitely many selected scales.}

{For $B(y_h,\varepsilon)\in\mathcal G_\varepsilon$, set}
\[
{
\widehat B_h:=\widehat B_{y_h,\bar\kappa,\varepsilon}\subset\subset\Omega,\qquad
E_h:=\{x\in E:\ x\in \widehat B_h,\ (1+\varepsilon)x\in \widehat B_h\}.
}
\]
{By Lemma~\ref{lem:flat-cover}, applied with $r=\varepsilon$ (so that the present set $E$ is $E_{\kappa,r}$),}
\[
{
E\setminus\bigcup_{B(y_h,\varepsilon)\in\mathcal G_\varepsilon}E_h
\subset
\bigcup_{B(y_h,\varepsilon)\in\mathcal R_\varepsilon}B(y_h,C_\kappa\varepsilon).
}
\]
		{For $x\in E_h$, letting $(\nabla u)_{\widehat B_h}:=\frac{1}{|\widehat B_h|}\int_{\widehat B_h}\nabla u$,
		we have}
\[
|\nabla u((1+\varepsilon)x)-\nabla u(x)|^{1+\gamma}
\le C_\gamma\Bigl(|\nabla u((1+\varepsilon)x)-(\nabla u)_{\widehat B_h}|^{1+\gamma}
+|\nabla u(x)-(\nabla u)_{\widehat B_h}|^{1+\gamma}\Bigr).
\]
Integrating over $E_h$ and using the change of variables $z=(1+\varepsilon)x$ in the first term
(note that $z\in\widehat B_h$ for $x\in E_h$) yields
\begin{equation}\label{7.2}
\int_{E_h}\frac{|\nabla u((1+\varepsilon)x)-\nabla u(x)|^{1+\gamma}}{2\varepsilon}\,dx
\le C\,\varepsilon^{-1}\int_{\widehat B_h}|\nabla u-(\nabla u)_{\widehat B_h}|^{1+\gamma}\,dx.
\end{equation}
By Sobolev--Poincar\'e applied to $\nabla u$ on $\widehat B_h$, and choosing
$\gamma>0$ so small that $1+\gamma\le \frac{n}{n-1}$, we get
\begin{equation}\label{7.3}
\int_{\widehat B_h}|\nabla u-(\nabla u)_{\widehat B_h}|^{1+\gamma}\,dx
\le {C(n,\gamma)\,(\bar\kappa^{-1}\varepsilon)^{1+\gamma}
\left(\int_{\widehat B_h}|D^2u|\,dx\right)^{1+\gamma}(\bar\kappa^{-1}\varepsilon)^{-n\gamma}.}
\end{equation}
Using \eqref{7.1} with $r=\varepsilon$ (since $y_h\in G_{\kappa,\eta}$ and $\varepsilon\le r_{\kappa,\eta}$),
\[
{
\int_{\widehat B_h}|D^2u|\,dx \le \delta_\kappa\,(\bar\kappa^{-1}\varepsilon)^{n-1},
}
\]
and plugging this into \eqref{7.3} and then \eqref{7.2} gives
\[
{
\int_{E_h}\frac{|\nabla u((1+\varepsilon)x)-\nabla u(x)|^{1+\gamma}}{2\varepsilon}\,dx
\le C(n,\gamma)\,\delta_\kappa^{1+\gamma}\,\bar\kappa^{-n}\,\varepsilon^{n-1}.
}
\]
{Summing over $B(y_h,\varepsilon)\in\mathcal G_\varepsilon$} and using the bounded overlap together with
the finiteness of the upper Minkowski content (so that the number of such balls is
{$\lesssim \varepsilon^{-(n-1)}$; see Footnote 4}) and \eqref{eq:deltakappa}, we obtain
\begin{equation}\label{2.8}
\sum_{B(y_h,\varepsilon)\in\mathcal G_\varepsilon}
\int_{E_h}\frac{|\nabla u((1+\varepsilon)x)-\nabla u(x)|^{1+\gamma}}{2\varepsilon}\,dx
\le {C(M,n,\mathrm{diam}(\Omega))\,\kappa^\gamma.}
\end{equation}

\smallskip
\noindent\emph{Bad part and choice of scales.}
{By Lemma~\ref{lem:flat-cover} and Proposition~\ref{prop:ADR}(\ref{item:badcount}), the remaining portion of $E$ is contained in
the $C_\kappa\varepsilon$-enlargement of at most}
\[
{
C(n,A_1,A_2)\,\varepsilon^{-(n-1)}\Bigl(|\log\varepsilon|^{-1}\eta^{-2}+\theta(\eta)\Bigr)
}
\]
{bad balls. We choose the gauge $\theta(\eta)=o_\eta(1)$, hence}
\begin{equation}\label{2.9}
{
\Bigg|E\setminus\bigcup_{B(y_h,\varepsilon)\in\mathcal G_\varepsilon}E_h\Bigg|
\le C(n,A_1,A_2,\kappa)\,\varepsilon
\Bigl(|\log\varepsilon|^{-1}\eta^{-2}+o_\eta(1)\Bigr).
}
\end{equation}
{Since $|\nabla u|\le L$, the contribution of this bad set to $I^+_{2,\varepsilon}$ is bounded by}
$$
{
C(n,A_1,A_2,L,\kappa)\,\bigl(|\log\varepsilon|^{-1}\eta^{-2}+o_\eta(1)\bigr).
}
$$
{Combining \eqref{2.7}, \eqref{2.8}, and \eqref{2.9}, we obtain the following estimate on $I^+_{2,\varepsilon}$
(for $\varepsilon=r_j$)}
$$
{I_{2, r_j}^+\le O(\kappa) +  C\,\kappa^\gamma
 + C(n,A_1,A_2,L,\kappa)\bigl(|\log r_j|^{-1}\eta^{-2}+o_\eta(1)\bigr),}
$$
{where the constant $C$ in the second term is independent of $\kappa,\eta$ and $j$.}


\medskip
\noindent\textit{Step 8: Conclusion.}
Fix $\kappa\in(0,1)$ and {then choose $0<\eta\le\eta_0(n,\kappa)$}. Let $r_{\kappa,\eta}>0$ be the scale produced in
Step~7 (so that the estimate \eqref{7.1} is valid for all
$0<\varepsilon\le r_{\kappa,\eta}$), {and let $c_\kappa$ be the geometric constant from Lemma~\ref{lem:flat-cover}}.
Let $\{r_j\}_{j\ge 1}$ be the sequence given in Step~7 by Proposition~\ref{prop:ADR} {with auxiliary set $G_{\kappa,\eta}$}.
Since $r_j\downarrow 0$, by discarding finitely many indices we may assume that
{$r_j\le c_\kappa r_{\kappa,\eta}$} for all $j$.

For $\varepsilon=r_j$, combining the bounds in Steps~5--7 yields
\[
I^+_{r_j}\le C\,\kappa^\gamma + O(\kappa)
+ {C(n,A_1,A_2,L,\kappa)}\Bigl(|\log r_j|^{-1}\eta^{-2}+o_\eta(1)\Bigr),
\]
and the same estimate holds for $I^-_{r_j}$. Hence, for $\varepsilon=r_j$,
\[
I^+_{r_j}+I^-_{r_j}\le C\,\kappa^\gamma + O(\kappa)
+ {C(n,A_1,A_2,L,\kappa)}\Bigl(|\log r_j|^{-1}\eta^{-2}+o_\eta(1)\Bigr).
\]
Letting $j\to\infty$ (so that $r_j\downarrow 0$ and $|\log r_j|^{-1}\to 0$), we obtain
\[
\limsup_{j\to\infty}\,\bigl(I^+_{r_j}+I^-_{r_j}\bigr)
\le C\,\kappa^\gamma + O(\kappa) + {C(n,A_1,A_2,L,\kappa)}\,o_\eta(1).
\]
Finally, we first let $\eta\downarrow 0$ (so that $o_\eta(1)\to 0$), and then let
$\kappa\downarrow 0$, to conclude
\[
\limsup_{j\to\infty}\, \bigl(I^+_{r_j}+I^-_{r_j}\bigr)=0.
\]
Since $I^+_\varepsilon+I^-_\varepsilon\ge 0$, this implies
\[
\liminf_{\varepsilon\to 0}\bigl(I^+_\varepsilon+I^-_\varepsilon\bigr)=0.
\]
As explained in Step~1, this yields \eqref{change derivative} and concludes the proof.
\end{proof}

\end{document}
